\section{Introduction}

Serious-game analytics depend on the chain of reasoning that links low-level events to higher-level claims about performance, learning, or adaptation. Telemetry design therefore cannot be treated as a secondary engineering concern. Chung \parencite{chung2015telemetry} argued that game telemetry should record behavior at a usable grain and preserve enough context to support valid inference. In health-oriented serious games, metadata and standardized descriptions are similarly important for interpretability and reuse \parencite{gobel2017metadata}.

The two cited approaches address different levels of description. Chung's telemetry guidance concerns the relationship between recorded player behavior, contextual information, and defensible performance inference \parencite{chung2015telemetry}. The health-game metadata model of G\"{o}bel and Maddison describes games as discoverable and comparable resources, including their intended users and purposes \parencite{gobel2017metadata}. Neither citation alone establishes a standardized, executable contract for a particular challenge--regulation--recovery episode. Our contribution is narrower: an event-level contract with cross-field validity rules and replay provenance. This comparison motivates the scope of the proposed artifact; it does not establish novelty relative to every existing event specification or prove construct validity.

Adaptive interventions create an additional problem. The system does not merely observe behavior; it changes task difficulty, sensory load, or support in response to prior events. If the event log omits the eligibility state, controller version, intervention assignment, timing basis, or reasons for invalidity, a later analyst may be unable to reconstruct why an adaptation occurred. This becomes especially important when the outcome itself is temporal, such as recovery after a controlled challenge.

We define a challenge--regulation--recovery event as a versioned scientific object rather than a collection of application logs. The objective is not to propose a universal psychological definition of recovery. Instead, the schema provides a reproducible representation of a study-specific recovery process once its operational definition has been frozen. The design is motivated by ARC-ADHD but is intentionally separable from therapeutic claims.

The paper makes four contributions: explicit event states and invariants across challenge and recovery; separation of missingness, technical invalidity, and non-recovery; preservation of controller and measurement context for replay; and machine-readable validation so conformance can be checked independently of the original application.
